\documentclass[11pt,a4paper]{article}
\usepackage[T1]{fontenc}
\usepackage[utf8]{inputenc}
\usepackage{lmodern}
\usepackage[margin=26mm]{geometry}
\usepackage{amsmath,amssymb,amsthm}
\usepackage{microtype}
\usepackage{hyperref}
\hypersetup{colorlinks=true,linkcolor=blue,citecolor=blue,urlcolor=blue,
 pdftitle={Global attraction of the positive equilibrium in an irreversible reaction-diffusion system},
 pdfauthor={Kaiwen Jin}}
\newtheorem{theorem}{Theorem}[section]
\newtheorem{lemma}[theorem]{Lemma}
\title{Global attraction of the positive equilibrium\\
in an irreversible reaction--diffusion system}
\author{Kaiwen Jin\thanks{\href{mailto:yui.ui.siki@gmail.com}{yui.ui.siki@gmail.com}}}
\date{15 October 2026}
\begin{document}
\maketitle
\begin{abstract}
An irreversible reaction--diffusion system can have both a positive
equilibrium and a stationary state at which the catalyst is absent.
For the three-species system studied by Nguyen and Tang, we prove that
every strictly positive compatible initial triple converges uniformly
to the positive equilibrium whenever the conserved masses satisfy
\(M_2\le M_1<2M_2\). The result holds on bounded connected smooth domains
in every dimension and for arbitrary positive diffusion coefficients.
Quadratic dissipation first removes the spatial oscillations of two
species. The total catalyst mass then satisfies a logistic equation
with a relative error tending to zero, which identifies the positive
limit and closes the strictly positive form of the global-attraction question.
\end{abstract}

\section{Introduction}\label{sec:intro}
We consider the reaction--diffusion system
\begin{equation}\label{eq:system}
 \begin{aligned}
 a_t-d_1\Delta a&=b(c-a),\\
 b_t-d_2\Delta b&=b(c-a),\\
 c_t-d_3\Delta c&=-b(c-a)
 \end{aligned}
 \qquad\text{in }\Omega,
\end{equation}
with homogeneous Neumann boundary conditions. Here \(\Omega\subset\mathbb R^n\)
is bounded, connected and smooth, \(|\Omega|=1\), and \(d_1,d_2,d_3>0\).
The concentrations \(a,b,c\) represent the irreversible chemical reaction
studied in \cite{NT}, in which the presence of \(b\) enables the reaction
between \(a\) and \(c\). Diffusion and the reaction preserve
\begin{equation}\label{eq:masses}
 \int_\Omega(a(t)+c(t))\,dx=M_1,\qquad
 \int_\Omega(b(t)+c(t))\,dx=M_2,
\end{equation}
where \(M_1=\int_\Omega(a_0+c_0)\,dx\) and
\(M_2=\int_\Omega(b_0+c_0)\,dx\).
For nonnegative bounded data the global classical solution is uniformly
bounded; this is established in \cite[Theorem~2.1]{NT}.
The remaining issue here is its long-time selection of equilibrium.

When \(M_2\le M_1<2M_2\), the conserved masses admit both the positive equilibrium
\[
 \left(\frac{M_1}{2},\,M_2-\frac{M_1}{2},\,\frac{M_1}{2}\right)
\]
and the boundary equilibrium \((M_1-M_2,0,M_2)\).
Nguyen and Tang prove local exponential stability of the positive
equilibrium in this regime \cite[Theorem~2.7]{NT}, and formulate global
attraction as a conjecture after their Table~1.
Their global positive-convergence result in the different regime
\(M_1<M_2\) uses the mass lower bound
\(\int_\Omega b\ge M_2-M_1>0\) \cite[Theorem~2.4]{NT}.
That bound is unavailable in the coexistence regime, so it does not
settle whether every positive initial state selects the positive equilibrium.

We prove this selection for strictly positive initial concentrations.
The fixed version of \href{https://github.com/MColbrook/AIM/blob/a153f469ee0ea64e876b2ad97c79d252c19a3e22/problems/416-catalytic-reaction-diffusion-positive-attractor.md}{AIM Problem~416}
assumes the same strictly positive \(C^2(\overline\Omega)\) data,
Neumann compatibility, and unique global bounded classical solution used here.
This restriction excludes the stationary boundary equilibrium:
the printed nonnegative-data wording of the source conjecture permits
\(b_0=0\), for which that equilibrium cannot evolve to a positive state.

\begin{theorem}\label{thm:main}
Let \(n\ge1\), let \(\Omega\subset\mathbb R^n\) be a bounded connected
smooth domain of volume one, and let \(d_1,d_2,d_3>0\).
Assume \(a_0,b_0,c_0\in C^2(\overline\Omega)\) are strictly positive
and satisfy the homogeneous Neumann compatibility conditions.
Let \((a,b,c)\) be the unique global bounded classical solution
of \eqref{eq:system}. If \(M_2\le M_1<2M_2\), then
\[
 \lim_{t\to\infty}
 \left[\left\|a(t)-\frac{M_1}{2}\right\|_\infty+
 \left\|b(t)-\left(M_2-\frac{M_1}{2}\right)\right\|_\infty+
 \left\|c(t)-\frac{M_1}{2}\right\|_\infty\right]=0.
\]
\end{theorem}
The masses are those in \eqref{eq:masses}, and all uniform norms are over
\(\overline\Omega\), equivalently essential uniform norms over \(\Omega\).
The diffusion coefficients may be unequal, and no smallness of the
initial deviation from equilibrium is required. The statement is
convergence for each admitted solution, with no common rate over all data.

The quadratic energy of \(a,c\), already used in the proof of
\cite[Theorem~2.7]{NT}, dissipates their spatial gradients even when the
catalyst is small. Compactness of bounded forced heat orbits turns this
dissipation into uniform spatial homogenization. The conserved masses
then reduce the catalyst's total mass to a positive logistic comparison
with an error controlled by those spatial oscillations.
Its limit identifies the means of \(a,c\); a final heat estimate removes
the remaining spatial variation of \(b\).
Section~\ref{sec:heat} gives the two heat lemmas, and
Section~\ref{sec:proof} proves the attraction theorem.
Throughout, \(\bar f=\int_\Omega f\,dx\) is the spatial average.

\section{Two heat equation lemmas}\label{sec:heat}
For a continuous real function \(f\) on \(\overline\Omega\), write
\[
 \bar f=\int_\Omega f\,dx,\qquad Qf=f-\bar f,\qquad
 X=C(\overline\Omega).
\]
The volume normalization makes \(\bar f\) the spatial average.
We use the pointwise uniform norm on \(X\) and the usual \(L^2(\Omega)\)
norm. The Neumann heat semigroup for \(u_t=d\Delta u\) is denoted by
\(S_d(t)\), where \(d>0\). Its constant mode is preserved, so \(Q\)
commutes with \(S_d(t)\).

We recall the precise heat properties used below. The Neumann realization
of \(-\Delta\) has a complete real orthonormal eigenbasis in \(L^2(\Omega)\)
and a smooth positive-time kernel; see \cite[Chapter~4, Section~8,
Proposition~4.56 and Theorem~4.58]{Baudoin}. Green's identity
\cite[Proposition~4.49]{Baudoin} shows that a zero eigenfunction has zero
gradient. Connectedness therefore leaves only the constants in the zero
eigenspace. We index the eigenvalues by
\[
 0=\lambda_0<\lambda_1\le\lambda_2\le\cdots,\qquad \psi_0=1.
\]
The multiplier of \(S_d(t)\) on \(\psi_j\) is \(e^{-d\lambda_jt}\).
For \(C^2\) functions \(g\) with \(\partial_\nu g=0\), Green's identity
gives
\begin{equation}\label{eq:poincare}
 \|\nabla g\|_2^2
 =\sum_{j\ge1}\lambda_j|\langle g,\psi_j\rangle|^2
 \ge\lambda_1\|Qg\|_2^2 .
\end{equation}
The coefficient inner product converges because \(g,\Delta g\in L^2\).

The heat flow preserves nonnegativity. One may test its equation by a
smooth convex approximation of the squared negative part, then use
strong \(L^2\) continuity at time zero. Its continuous kernel \(k_d(t,x,y)\)
is consequently nonnegative: a negative kernel value would give a negative
heat value for a nonnegative test function supported nearby.
The constant mode gives
\[
 \int_\Omega k_d(t,x,y)\,dy=1,\qquad
 \|S_d(t)f\|_\infty\le\|f\|_\infty.
\]
At each fixed \(\tau>0\), continuity of the kernel on the compact closure
product also gives
\begin{equation}\label{eq:smoothing}
 \|S_d(\tau)f\|_\infty\le H_{d,\tau}\|f\|_2,\qquad
 H_{d,\tau}:=\sup_{x\in\overline\Omega}
       \left(\int_\Omega k_d(\tau,x,y)^2\,dy\right)^{1/2}<\infty .
\end{equation}
Bounded input sets in \(X\) have uniformly bounded, equicontinuous images
under \(S_d(\tau)\). A finite spatial net and finite grids for the image
values on that net give finite covers in the uniform norm. Thus
\(S_d(\tau)\) maps bounded sets into relatively compact sets in \(X\).

We will also use, for every \(C^2\) Neumann-compatible \(g\),
\begin{equation}\label{eq:shorttime}
 S_d(h)g-g=d\int_0^h S_d(r)\Delta g\,dr,\qquad
 \|S_d(h)g-g\|_\infty\le dh\|\Delta g\|_\infty .
\end{equation}
For the identity, Green's identity gives
\[
 \langle\Delta g,\psi_j\rangle=-\lambda_j\langle g,\psi_j\rangle,
\]
since both functions satisfy the Neumann condition. Integrating the
multiplier \(e^{-d\lambda_jr}\) proves \eqref{eq:shorttime} on each
spectral coefficient, including the zero coefficient. The integral
exists in \(X\): contraction bounds every tail over an interval of
length \(\varepsilon\) by \(\varepsilon\|\Delta g\|_\infty\), so the
truncated integrals are Cauchy in \(X\). Completeness of the eigenbasis
identifies the two sides in \(L^2\); their continuity gives the identity
in \(X\). Contraction then gives the displayed norm estimate.
These arguments use fixed positive heat times and are valid in every dimension.

\begin{lemma}\label{lem:compact}
Let \(d>0\), and let \(u_t-d\Delta u=f\) be a classical solution with
homogeneous Neumann conditions for positive times. Suppose that \(f\)
is continuous in time with values in \(X\), and that
\[
 \sup_{t\ge1}\|u(t)\|_\infty<\infty,\qquad
 F:=\sup_{t\ge1}\|f(t)\|_\infty<\infty.
\]
Then \(\{u(t):t\ge2\}\) has compact closure in \(X\), and \(u\) is
uniformly continuous from \([2,\infty)\) into \(X\).
\end{lemma}
\begin{proof}
The forced heat formula is
\begin{equation}\label{eq:duhamel}
 u(t)=S_d(t-s)u(s)+\int_s^tS_d(t-r)f(r)\,dr,\qquad t\ge s\ge1.
\end{equation}
It follows by solving the scalar equation on each Neumann eigenfunction
coefficient. The integrand is continuous in \(X\) away from its upper
endpoint. The improper integral converges in \(X\), because every
endpoint tail of length \(\varepsilon\) has norm at most \(F\varepsilon\).
Completeness of the eigenbasis then gives the stated identity.

For \(t\ge2\) and \(0<\varepsilon<1\), split the last unit time interval:
\begin{align*}
 u(t)={}&S_d(1)u(t-1)
 +S_d(\varepsilon)\int_\varepsilon^1
           S_d(r-\varepsilon)f(t-r)\,dr\\
 &+\int_0^\varepsilon S_d(r)f(t-r)\,dr .
\end{align*}
The first term ranges over a relatively compact set. The input of the
middle heat operator is bounded by \(F\), so the second term also ranges
over a relatively compact set. The last term has norm at most
\(F\varepsilon\). Taking \(\varepsilon\) small and covering the first
two sets proves total boundedness of the orbit. Its closure is compact
because \(X\) is complete.

For \(\eta>0\), choose a finite \(\eta\)-net with centers
\(g_j=u(t_j)\) from the actual orbit, \(t_j\ge2\). Each \(g_j\) is
\(C^2\) and Neumann-compatible. Contraction and \eqref{eq:shorttime},
followed by \eqref{eq:duhamel}, give
\[
 \|u(t+h)-u(t)\|_\infty
 \le2\eta+h\left(d\max_j\|\Delta g_j\|_\infty+F\right),\qquad t\ge2 .
\]
The maximum is finite because there are finitely many centers. First
choose \(\eta\) small, then \(h\) small. This proves uniform continuity;
it requires no bound on \(\Delta u(t)\) over the entire orbit.
\end{proof}

\begin{lemma}\label{lem:forced}
Let \(d>0\), and let \(w\) be a classical Neumann solution of
\(w_t-d\Delta w=Qf\), with \(\bar w(t)=0\). Suppose that \(f(t)\in X\)
is continuous, bounded for \(t\ge1\), and
\(\|f(t)\|_\infty\to0\). Then \(\|w(t)\|_\infty\to0\).
\end{lemma}
\begin{proof}
Set \(\gamma=d\lambda_1>0\). The spectral expansion and
\eqref{eq:smoothing} imply, for \(r\ge1\),
\[
 \|S_d(r)Qf\|_\infty
 \le H_{d,1}e^{-\gamma(r-1)}\|Qf\|_2
 \le H_{d,1}e^{-\gamma(r-1)}\|f\|_\infty.
\]
For \(0\le r\le1\), contraction gives \(2\|f\|_\infty\).
Thus, for a finite constant \(C_d\) depending on \(d,\Omega\),
\begin{equation}\label{eq:decay}
 \|S_d(r)Qf\|_\infty
 \le C_de^{-\gamma r}\|f\|_\infty,\qquad r\ge0 .
\end{equation}
Apply \eqref{eq:duhamel} from any fixed \(T\ge2\):
\[
 \|w(t)\|_\infty
 \le C_de^{-\gamma(t-T)}\|w(T)\|_\infty+
 C_d\int_T^t e^{-\gamma(t-s)}\|f(s)\|_\infty\,ds .
\]
For a fixed \(S\ge T\), the integral over \([T,S]\) tends to zero as
\(t\to\infty\), while the part over \([S,t]\) is at most
\(\gamma^{-1}\sup_{s\ge S}\|f(s)\|_\infty\).
Taking \(t\to\infty\) and then \(S\to\infty\) proves the lemma.
\end{proof}

\section{Proof of the attraction theorem}\label{sec:proof}
\begin{proof}[Proof of Theorem~\ref{thm:main}]
Fix a uniform bound \(K\ge1\) for the absolute values of \(a,b,c\).
First, all three components remain nonnegative. To see this directly,
write \(U=(a,b,c)\), \(p_i=U_i^-\), and \(V=U^+\), and denote
the reaction vector by \(F(U)=(b(c-a),b(c-a),-b(c-a))\).
On the set \(p_i>0\), the \(i\)-th reaction component satisfies
\(F_i(V)\ge0\), whereas
\[
 |b(c-a)-V_b(V_c-V_a)|\le2K(p_a+p_b+p_c).
\]
Testing the three equations by the negative parts, justified by a smooth
convex approximation, gives
\[
 P_-(t)\le12K\int_0^tP_-(s)\,ds,\qquad
 P_-(t):=\int_\Omega(p_a^2+p_b^2+p_c^2)\,dx .
\]
Here \(P_-(0)=0\). For \(J(t)=\int_0^tP_-(s)\,ds\),
\((e^{-12Kt}J(t))'\le0\), and \(J\ge0\); hence \(P_-=0\).

Integrating the equations preserves the two masses in
\eqref{eq:masses}. In particular, with \(B(t)=\bar b(t)\),
\begin{equation}\label{eq:means}
 \bar a=M_1-M_2+B,\qquad
 \bar c=M_2-B,\qquad B'=\int_\Omega b(c-a)\,dx .
\end{equation}
Since \(c-a\ge-K\),
\[
 B(t)\ge B(0)e^{-Kt}>0\qquad\text{for every finite }t\ge0 .
\]
Only this finite-time nonvanishing will be needed when normalizing by \(B\).

Consider the quadratic energy of \(a,c\):
\[
 E(t)=\frac12\int_\Omega(a(t)^2+c(t)^2)\,dx.
\]
The Neumann conditions remove the boundary terms, and
\((a-c)b(c-a)=-b(c-a)^2\) gives the exact dissipation identity
\begin{equation}\label{eq:energy}
 E'=-d_1\|\nabla a\|_2^2-d_3\|\nabla c\|_2^2
       -\int_\Omega b(c-a)^2\,dx\le0 .
\end{equation}
Since \(E\ge0\), the two gradient squares are time-integrable.
By \eqref{eq:poincare},
\[
 z(t):=\|Qa(t)\|_2^2+\|Qc(t)\|_2^2,\qquad
 \int_2^\infty z(t)\,dt<\infty .
\]
The forcing \(b(c-a)\) is uniformly bounded and continuous in \(X\).
Lemma~\ref{lem:compact} applies separately to \(a,c\) with their own
diffusivities, and gives uniform continuity in time in \(X\).
The quadratic functional defining \(z\) is Lipschitz on bounded subsets
of \(X\); for example,
\[
 \big|\|Qg\|_2^2-\|Qh\|_2^2\big|
 \le8K\|g-h\|_\infty\quad
 \text{if }\|g\|_\infty,\|h\|_\infty\le K.
\]
Thus \(z\) is uniformly continuous. A nonnegative uniformly continuous
function with finite integral tends to zero: otherwise a separated
sequence of times would give disjoint intervals of a common positive
length on which the function stays bounded below. Therefore \(z(t)\to0\).

The same lemma makes the \(a,c\) orbits relatively compact in \(X\).
It follows that
\begin{equation}\label{eq:oscillations}
 h(t):=\|a(t)-\bar a(t)\|_\infty+
       \|c(t)-\bar c(t)\|_\infty\longrightarrow0 .
\end{equation}
Indeed, a violating sequence has a common uniformly convergent
subsequence for \(a,c\). The limits have zero variance because
\(z\to0\), and hence are constant, contradicting that violation.

It remains to identify these changing averages. Put
\(\delta=2M_2-M_1>0\). Equation~\eqref{eq:means} gives
\(\bar c-\bar a=\delta-2B\), so the exact equation for \(B\) is
\begin{equation}\label{eq:meanode}
 B'=B(\delta-2B+\rho(t)),\qquad
 \rho(t)=\frac1{B(t)}\int_\Omega
          b\big[(c-\bar c)-(a-\bar a)\big]\,dx .
\end{equation}
Because \(b\ge0\) and \(\int_\Omega b=B>0\),
\begin{equation}\label{eq:rhobound}
 |\rho(t)|\le h(t)\longrightarrow0 .
\end{equation}
The weight \(b/B\) has total mass one, so this estimate remains valid
even when \(B\) is small.

For \(0<\eta<\delta\), choose \(T\ge2\) with
\(|\rho(t)|\le\eta\) for every \(t\ge T\). Then
\[
 B(\delta-\eta-2B)\le B'\le B(\delta+\eta-2B).
\]
For \(k_\pm=\delta\pm\eta\), define
\[
 y_\pm(t)=\frac{k_\pm/2}
 {1+\big(k_\pm/(2B(T))-1\big)e^{-k_\pm(t-T)}} .
\]
The denominators are positive, since \(k_\pm>0\) and \(B(T)>0\).
These functions solve \(y_\pm'=y_\pm(k_\pm-2y_\pm)\) and start at
\(B(T)\). The difference \(D=B-y_-\) satisfies
\[
 D'\ge\big(k_--2(B+y_-)\big)D,\qquad D(T)=0 ,
\]
so multiplication by the corresponding positive integrating factor
gives \(B\ge y_-\). Applying the same calculation to \(y_+-B\)
gives \(B\le y_+\). Hence
\[
 \frac{\delta-\eta}{2}\le\liminf_{t\to\infty}B(t)
 \le\limsup_{t\to\infty}B(t)\le\frac{\delta+\eta}{2}.
\]
Letting \(\eta\downarrow0\) proves
\[
 B(t)\longrightarrow M_2-M_1/2,\qquad
 \bar a(t),\bar c(t)\longrightarrow M_1/2 .
\]
Together with \eqref{eq:oscillations}, this identifies the uniform
limits of \(a,c\). In particular
\[
 \|b(t)(c(t)-a(t))\|_\infty
 \le K\big(\|c(t)-M_1/2\|_\infty+\|a(t)-M_1/2\|_\infty\big)\to0 .
\]
Finally \(w=b-B\) has zero average and satisfies
\[
 w_t-d_2\Delta w=Q\big(b(c-a)\big).
\]
Lemma~\ref{lem:forced} gives \(\|b(t)-B(t)\|_\infty\to0\).
The already determined limit of \(B\) proves the required limit of \(b\),
and completes the convergence of all three components.
\end{proof}

\section*{AI use}
OpenAI Codex assisted with source retrieval, development and checking of the
mathematical argument, and preparation of this manuscript. Two further
evaluations in fresh Codex contexts provided reader and referee feedback
on one frozen draft; they shared the same provider and workspace with the
drafting process. The manuscript was then revised. Independent mathematical
review is welcome.

\begin{thebibliography}{9}
\bibitem{NT}
T.~L. Nguyen and B.~Q. Tang,
\emph{Stability analysis of irreversible chemical reaction-diffusion
systems with boundary equilibria},
Zeitschrift f\"ur angewandte Mathematik und Physik
\textbf{77} (2026), article 199, 30 pp.
\href{https://doi.org/10.1007/s00033-026-02847-0}
{doi:10.1007/s00033-026-02847-0}.

\bibitem{Baudoin}
F. Baudoin,
\emph{Diffusion processes and Stochastic calculus},
author version dated 13 February 2019,
Chapter~4, Proposition~4.49 and Section~8, pp.~131, 135--136.
\url{https://fabricebaudoin.blog/wp-content/uploads/2019/12/book_emsstyle.pdf}.
\end{thebibliography}
\end{document}
